\chapter{Sphere Factors and Connected-Sum Reduction}
\label{ch:sphere-reduction}

Finite extinction turns the evolution into a finite topological assembly.
The remaining task is to recognize its factors and then its target.  Simple
connectedness does the first job: each factor group is a retract of the
initial fundamental group, a compact sphere bundle over a circle cannot be
simply connected, and a simply connected positive spaceform is a sphere.
The second job requires control of the actual attaching balls and collars.
We construct a sphere map for a binary operation and propagate a finite
disjoint-union invariant through the assembly.

Throughout this chapter a smooth three-manifold is Hausdorff and second
countable, has no boundary, and is modelled on $\mathbb R^3$.
It may be disconnected unless connectedness is specified.  The original
manifold is denoted $M$, the chosen flow has slices $M_t$, and
$S^3=\{x\in\mathbb R^4:\|x\|=1\}$ has its standard smooth structure.
The Borel measurable structure is used when applying the flow theorems;
none of the topological arguments below depends on an additional measure.

\section{The exact reconstruction interface}
\label{sec:sphere-reconstruction-interface}

Fix a normalized initial metric on a connected $M$, one controlled surgery
flow, its initial diffeomorphism $e_0:M\to M_0$, and a finite extinction
time $T_{\rm ext}$.  The finite-history construction supplies local
decompositions at precisely the events
\[
 \mathcal E=\Sigma\cap[0,T_{\rm ext}],
\]
where $\Sigma$ is this flow's surgery-time set.  At $e\in\mathcal E$ write
$Q_{e,i}$, $0\le i<p_e$, for the nonempty compact connected local pieces.
Each is labelled survivor, sphere bundle, or positive spaceform.  A survivor
is identified with a component of the outgoing slice.  The other two kinds
are recorded even when the outgoing slice is empty.

\begin{definition}[Exact factor family]
\label{def:exact-factor-ledger}
The discarded-factor index set is
\[
 \mathcal J=\{(e,i):e\in\mathcal E,\ 0\le i<p_e,\
                       Q_{e,i}\text{ is not labelled survivor}\}.
\]
Choose a bijection $\sigma:\{0,\ldots,r-1\}\to\mathcal J$ and put
$P_j=Q_{\sigma(j)}$.  An exact reconstruction consists of this family,
the specified equalities with the event pieces, and a finite smooth
connected-sum assembly of $(P_j)_{j<r}$ into $M_0$.
\lean{PoincareMT.M72ReconstructionConclusion}
\uses{prop:connected-sum-reconstruction}
\source{MT2007}
\dcref{Proposition 15.3 and Corollary 15.4, pp. 357--359}
\end{definition}

A smooth disjoint-union presentation of $X$ by $(P_j)$ means a finite
partition $X=\coprod_j U_j$ into clopen sets and diffeomorphisms $P_j\to U_j$.
A binary connected-sum operation has two side manifolds $A,B$, disjoint
open images of their punctured sides in a new manifold, and a smooth
sphere collar joining them.  Precise ball and collar data are given in
Section~\ref{sec:binary-sphere-gluing}.  A finite assembly starts with a
smooth disjoint-union presentation of the given factors and performs
finitely many such binary operations.  Side manifolds may themselves be
disconnected.  Thus the intermediate object is not asserted to be a single
connected sum of connected manifolds.

Proposition~\ref{prop:connected-sum-reconstruction} provides an assembly into the exact $M_0$,
which is nonempty and connected.  It also establishes that $\sigma$ is
bijective: every discarded event piece appears once.  Its successor
transports and reverse induction belong to that theorem.  Here we need its
geometric output, including the actual gluing data, rather than a second
choice of manifolds known only to have the same diffeomorphism types.

\section{Fundamental groups of assembly factors}
\label{sec:factor-fundamental-groups}

\begin{lemma}[Retractions through a finite assembly]
\label{lem:sphere-factor-retraction}
\uses{def:finite-smooth-assembly}
For every finite smooth connected-sum assembly $(P_j)_{j<r}\leadsto X$,
every $j<r$, and every $x\in P_j$, there are a point $y\in X$ and group
homomorphisms
\[
 \iota:\pi_1(P_j,x)\longrightarrow\pi_1(X,y),\qquad
 p:\pi_1(X,y)\longrightarrow\pi_1(P_j,x),\qquad p\iota=1.
\]
In particular, if $X$ is simply connected and $P_j$ is nonempty and connected, then
$P_j$ is simply connected.
\lean{PoincareMT.SmoothFiniteConnectedSumAssembly.piece_factor}
\lean{PoincareMT.SmoothFiniteConnectedSumAssembly.piece_simplyConnected}
\source{MT2007}
\dcref{Proposition 15.3, pp. 357--358}
\end{lemma}

\begin{proof}
For the initial disjoint union, a based loop and any homotopy of it stay
in the clopen region containing the basepoint.  Its fundamental group
is therefore the group of the corresponding factor.

Consider one binary operation, with punctured-side regions $U_A,U_B$ and
collar $c:S^2\times(-1,1)\to Y$.  The open cover
\[
 U_A,\qquad U_B\cup c(S^2\times(-1,1))
\]
has intersection $c(S^2\times(-1,0))$.  This intersection is simply
connected, since $S^2$ is simply connected and an interval is contractible.
Van Kampen's theorem \cite[Theorem 1.20, p.~43]{Hatcher} gives a homomorphism from the group of
$Y$, based in $U_A$, onto the group of $U_A$ whose composite with the
inclusion homomorphism is the identity.  It is the usual free-product
retraction: keep loops in the first set and send loops from the other set
to the identity.  The trivial intersection group makes these prescriptions
compatible.  If the side is disconnected, this is applied to the component
of the basepoint; components disjoint from the collar retain their groups.

Removing a collared closed three-ball does not change the corresponding
based fundamental group, up to a choice of basepoint in its complement.
One obtains this by filling the ball back in and applying van Kampen to
the exterior enlarged by a collar and the ball enlarged by a collar:
the ball is simply connected and the intersection retracts to $S^2$.
Transport across the punctured-side diffeomorphism and this group
isomorphism gives the displayed retraction for the original first side.
The positive half-collar gives the same conclusion for the second side.
More explicitly, for a radius-two ball chart $b$, use the open cover
$A\setminus b(\overline B_1)$ and $b(B_2)$; the overlap is the annulus
$b(\{1<|z|<2\})$.  If the original basepoint lies in the removed ball,
a path inside $b(B_2)$ joins it to $b(3q/2)$ for any $q\in S^2$.
Basepoint transport along this path gives the required isomorphism.
This proves the ball-removal input at every basepoint without assuming
that the side manifold is connected.  The same retractions are used for
survivors in Proposition~\ref{prop:filling-group-effects}.

Retractions compose.  If $p_1\iota_1=1$ and $p_2\iota_2=1$, then
\[
  (p_1p_2)(\iota_2\iota_1)=1.
\]
Induction through the finite sequence, starting with the disjoint-union
isomorphism, proves the assertion at every original basepoint.
If $\pi_1(X,y)$ is trivial, then for every $a\in\pi_1(P_j,x)$,
$a=p(\iota(a))=p(1)=1$.  Finally a connected smooth manifold is locally
path connected and hence path connected; trivial based fundamental groups
then give simple connectedness.
\end{proof}

If $M$ is simply connected, $e_0:M\to M_0$ transports simple connectedness
to $M_0$.  Apply the lemma to the exact assembly and use the connectedness
of each local event piece.  Every $P_j$ is simply connected.  Notice that
the argument produces the retractions from the assembly: merely having
a homomorphism into a trivial group would not suffice.

\section{Excluding bundles and recognizing spaceforms}
\label{sec:factor-classification}

\begin{definition}[The two discarded types]
\label{def:discarded-factor-types}
A sphere-bundle structure on $Q$ consists of a smooth continuous surjection
$\pi:Q\to S^1$ such that, for every $b\in S^1$, some open neighborhood
$V$ of $b$ has a smooth trivialization
$\pi^{-1}(V)\cong S^2\times V$ commuting with the second projection.
Both orientations of the monodromy are allowed.

A positive-spaceform structure on $Q$ consists of a smooth Riemannian
metric $g$, its Levi--Civita connection, compactness and connectedness of
$Q$, and a constant $c>0$ such that every sectional curvature equals $c$.
Equivalently, the equality holds for every orthonormal pair at every
point; the constant precedes the choices of point and plane.
\lean{PoincareMT.SurgerySphereBundle}
\lean{PoincareMT.SurgeryPositiveSpaceform}
\source{MT2007}
\dcref{Theorem 1.11(2), pp. 8--9}
\dcref{Proposition 15.3, pp. 357--358}
\end{definition}

\begin{lemma}[Compact bundle obstruction]
\label{lem:sphere-bundle-obstruction}
\uses{def:discarded-factor-types}
A compact smooth three-manifold admitting the sphere-bundle structure
of Definition~\ref{def:discarded-factor-types} is not simply connected.
\lean{PoincareMT.SurgerySphereBundle.not_simplyConnectedSpace}
\source{MT2007}
\dcref{Corollary 15.4(2), pp. 358--359}
\end{lemma}

\begin{proof}
Use the exponential covering $q:\mathbb R\to S^1$,
$q(t)=(\cos t,\sin t)$.  The bundle projection $\pi$ is open: a
trivialization carries an open subset of $\pi^{-1}(V)$ to an open subset
of $S^2\times V$, and projection onto $V$ is open.

Choose $x_0\in Q$ with $\pi(x_0)=q(0)$, using surjectivity.  If $Q$ were
simply connected, covering-space lifting
\cite[Proposition 1.33, pp.~61--62]{Hatcher}, together with its local path
connectedness, would give a continuous map $L:Q\to\mathbb R$ with
$L(x_0)=0$ and $qL=\pi$.  Its range $R=L(Q)$ is nonempty and compact,
hence closed in $\mathbb R$.

We verify the additional fact that $R$ is open.  At $t=L(x)$ choose an
open interval $I$ containing $t$ on which $q$ is injective and is a
homeomorphism onto an open arc.  Set $W=L^{-1}(I)$.  Then
\[
       V_t=I\cap q^{-1}(\pi(W))
\]
is an open neighborhood of $t$.  If $u\in V_t$, there is $z\in W$
with $\pi(z)=q(u)$.  Both $L(z)$ and $u$ lie in $I$, and
$q(L(z))=\pi(z)=q(u)$, so injectivity on $I$ gives $u=L(z)$.
Thus $V_t\subset R$.  Consequently $R$ is nonempty and clopen in the
connected line, so $R=\mathbb R$.  This contradicts its compactness.
\end{proof}

This excludes the twisted bundle as well as the product bundle, without
computing either fundamental group or assuming orientability.

\begin{proposition}[Simply connected positive spaceforms]
\label{prop:positive-spaceform-sphere}
\uses{def:discarded-factor-types}
A simply connected positive spaceform of dimension three is smoothly
diffeomorphic to $S^3$.
\lean{PoincareMT.SurgeryPositiveSpaceform.nonempty_diffeomorph_threeSphere}
\source{MT2007}
\dcref{Theorem 1.11(2), pp. 8--9}
\end{proposition}

\begin{proof}
Let $c>0$ be its curvature constant and put $g_1=cg$.
Constant rescaling leaves the Levi--Civita connection unchanged; for a
nondegenerate plane the numerator in sectional curvature scales by $c$
and its Gram determinant scales by $c^2$.  Therefore
$K_{g_1}=c^{-1}K_g=1$.  Compactness gives completeness.
The spherical-covering form of the constant-curvature theorem supplies
a surjective local diffeomorphism $q:S^3\to Q$.
This is the standard spaceform theorem
\cite[Theorem 1.11(2) and the following spaceform paragraph, pp.~8--9]{MT2007}, applied after the displayed
normalization, rather than a sphere-recognition hypothesis about $Q$.

A connected covering of a simply connected, locally path-connected
space has one sheet.  Indeed, lift the identity of the base with a chosen
value in the fiber to a section; uniqueness of lifts
\cite[Proposition 1.34, p.~62]{Hatcher} makes the section
composed with the projection the identity on the connected total space.
Thus $q$ is bijective.  Its local smooth inverses agree wherever their
domains overlap, so they form a global smooth inverse, as required.
\end{proof}

\begin{theorem}[Recognition of the exact factors]
\label{thm:exact-sphere-factors}
For every reconstruction of Definition~\ref{def:exact-factor-ledger}
with simply connected original manifold $M$, each $P_j$ is labelled
positive spaceform and admits a diffeomorphism $d_j:P_j\to S^3$.
\lean{PoincareMT.m73SphereFactors}
\lean{PoincareMT.m73SphereFactors_from_M72}
\source{MT2007}
\dcref{Proposition 15.3 and Corollary 15.4(2), pp. 357--359}
\uses{def:exact-factor-ledger,lem:sphere-factor-retraction,lem:sphere-bundle-obstruction,prop:positive-spaceform-sphere}
\end{theorem}

\begin{proof}
The retraction lemma makes each exact factor simply connected.
Its index excludes the survivor label.  The bundle label would supply
a sphere bundle on the compact factor, contradicting the bundle
obstruction.  Hence its label is positive spaceform, and the preceding
proposition supplies $d_j$.  Every use of compactness, connectedness,
and classification concerns $Q_{\sigma(j)}=P_j$, so the conclusion
retains the original indexing.
\end{proof}

\section{Binary smooth sphere gluing}
\label{sec:binary-sphere-gluing}

\begin{definition}[Collared connected-sum data]
\label{def:sphere-binary-data}
A surgery ball in $A$ is a smooth coordinate embedding
$b:B(0,2)\to A$ with smooth inverse on its open image.  Write
$D_b=b(\overline B(0,1))$.  Binary connected-sum data on $A,B,C$ consist
of balls $b_A,b_B$, disjoint open regions $V_A,V_B\subset C$,
diffeomorphisms
\[
 a:A\setminus D_{b_A}\to V_A,\qquad
 b:B\setminus D_{b_B}\to V_B,
\]
a diffeomorphism $\varphi:S^2\to S^2$, and a smooth collar embedding
$c:S^2\times(-1,1)\to C$ with smooth inverse on its open image, such that
\[
 \begin{aligned}
 c(q,s)&=a(b_A((1-s)q))&&(-1<s<0),\\
 c(q,s)&=b(b_B((1+s)\varphi(q)))&&(0<s<1).
 \end{aligned}
\]
The central sphere is disjoint from both regions, and
$C=V_A\cup V_B\cup c(S^2\times\{0\})$.
No orientation condition is imposed on $\varphi$.
\lean{PoincareMT.SmoothConnectedSumData}
\source{MT2007}
\dcref{Proposition 15.3 and Corollary 15.4, pp. 357--359}
\end{definition}

We use Lemmas~\ref{lem:neck-cap-prescribed-collar}
and~\ref{lem:neck-cap-sphere-isotopy} from the neck/cap chapter.
The collared Schoenflies input has the following precise form.
For every smooth injective immersion
$\psi:S^2\times(-1,1)\to\mathbb R^3$ and every $0<\delta<1$,
there are a sign $\eta\in\{-1,1\}$, a bounded connected open region
$\Omega$ disjoint from $\psi(S^2\times\{0\})$, a radius $R_*>0$,
a ball chart $F:B(0,R_*)\to\mathbb R^3$ with smooth inverse on its
image, a sphere diffeomorphism $\beta$, and a strictly increasing
function $r:[\delta,1)\to(0,R_*)$ such that
\[
 \begin{split}
 \psi(S^2\times\eta(-1,0))&\subset\Omega,\\
 F(B(0,R_*))&=\Omega\cup\psi(S^2\times\eta[0,1)),\\
 F(r(t)\beta(q))&=\psi(q,\eta t)
       \quad(q\in S^2,\ \delta\le t<1).
 \end{split}
\]
The complement of $\Omega\cup\psi(S^2\times\{0\})$ is connected.
No smoothness of $r$ is assumed: where needed it is recovered from
the inverse chart.  This is a stronger collar specification than the
bare assertion that an embedded sphere bounds a ball.
The sphere-isotopy input gives, for every diffeomorphism $f:S^2\to S^2$,
an orthogonal map $O$ and a jointly smooth family of diffeomorphisms
$I_t$, $t\in\mathbb R$, with $I_0=O|_{S^2}$ and $I_1=f$.
The orthogonal endpoint permits both orientation classes.
Their ordinary background is
\cite[Theorem 1.1 and sphere extension, pp.~1--5]{Hatcher3M2014}
and \cite[Theorems 5--6, pp.~624--626]{Smale1959};
the collar adjustment and pole-normalization arguments are developed in
Section~\ref{sec:neck-cap-prescribed-collars}.

\begin{lemma}[Reciprocal coordinates on a punctured sphere]
\label{lem:reciprocal-sphere-end}
\uses{lem:neck-cap-prescribed-collar,lem:neck-cap-sphere-isotopy}
Given $d:A\to S^3$ and a surgery ball $b$ in $A$, there exist a
diffeomorphism $E:A\setminus D_b\to\mathbb R^3$ and $0<\epsilon<1$
such that, for every $q\in S^2$ and $0<s<\epsilon$,
\[
             E(b((1+s)q))=s^{-1}q.
\]
\lean{PoincareMT.SurgeryBallEmbedding.exists_canonicalPuncturedSphereEnd}
\source{MT2007}
\dcref{Collared implementation of Corollary 15.4(2), pp. 358--359}
\end{lemma}

\begin{proof}
Let $e:A\setminus\{b(0)\}\to\mathbb R^3$ be $d$ followed by
stereographic projection having pole $d(b(0))$.  Put
$\psi(q,s)=e(b((1+s)q))$.  The change of parameter
\[
 v(t)=\frac{1+2t}{2+t},\qquad -1<t<1,
\]
is an increasing diffeomorphism of $(-1,1)$ onto itself.
The shifted collar $\widetilde\psi(q,t)=\psi(q,v(t))$ has its
central sphere at ball radius $3/2$, and reaches the original surgery
sphere at $t=-1/2$.  Apply the stated Schoenflies input to this collar
with $\delta=1/4$.

Its sign is $\eta=-1$: the negative half of the shifted collar is the
whole punctured ball of radius $3/2$ in $e$-coordinates, an unbounded
set because it approaches the stereographic pole.  It cannot lie in
the bounded region $\Omega$.  In fact $\Omega$ is precisely the image
of the exterior of the closed radius-$3/2$ ball.  To see the equality,
partition the complement of the central sphere into the punctured
inner-ball image and the outer image.  They are disjoint open sets.
Connectedness places $\Omega$ wholly in the outer set, since its
prescribed inner collar meets that set; connectedness of the other
Schoenflies side places it wholly in the inner set, since it meets
the negative half-collar.  The two containments force equality.
The chart-image identity therefore makes $F:B(0,R_*)\to\mathbb R^3$
surjective: its nonnegative oriented half-collar fills the punctured
inner ball and its boundary sphere.

Put $R=r(1/2)$.  The radial formula maps the sphere of radius $R$ to
the original surgery sphere.  It maps $B(0,R)$ exactly to
$e(A\setminus D_b)$.  Here is the side check for this smaller ball.
The chart images of $B(0,R)$ and $\{R<\|x\|<R_*\}$ are connected
and partition the complement of that sphere.  The first is bounded,
since $F$ is continuous on the compact closed radius-$R$ ball.
The original punctured inner-ball image is unbounded, so it meets
the annulus image.  Connectedness places that whole annulus image
in the inner side.  The nonempty exterior therefore lies in the
ball image, whose connectedness places it wholly in the exterior.
This proves the asserted equality without presupposing that the
original exterior is a ball.

The inverse chart on this exterior has the expression
\[
       b((1+s)q)\longmapsto \rho(s)\,\beta(q)
\]
near its boundary, where
\[
 \rho(s)=r\left(\frac{1-2s}{2-s}\right),\qquad |s|<1/8.
\]
The argument of $r$ lies in $(1/4,1)$, so this is within the prescribed
radial interval, and $\rho(0)=R$.  Strict monotonicity of $r$ makes
$\rho$ strictly decreasing.  For a fixed $q$, the norm of
$F^{-1}(\psi(q,s))$ equals $\rho(s)>0$, proving smoothness.
Differentiating
$F(\rho(s)\beta(q))=\psi(q,s)$ shows $\rho'(s)\ne0$, since the
right side has nonzero derivative by the immersion property of the
collar.  Thus $\rho'<0$, and the inverse function theorem gives its
smooth inverse near the boundary radius.

The sphere-isotopy input extends $\beta$ to a norm-preserving
diffeomorphism $K$ of $\mathbb R^3$: for $x=rq\ne0$ take
$K(rq)=rI_{\chi(r)}(q)$, with a smooth cutoff $\chi$ equal to zero
near the origin and one on an outer annulus.  Near zero this is the
linear orthogonal map, so there is no polar-coordinate singularity;
the same argument applies to its inverse.  Applying $K^{-1}$ removes
the angular term $\beta$.

It remains to replace $\rho(s)$ by $1/s$.  Let
$\tau=\rho^{-1}$ near $R$.  It is positive below $R$,
$\tau'<0$, and $\tau(r)\to0$ as $r\uparrow R$.
Consequently $h(r)=1/\tau(r)$ has positive derivative there and
diverges at $R$.  Extend this outer radial function to a smooth
increasing bijection $f:(0,R)\to(0,\infty)$ which is $kr$ near zero,
for some $k>0$, and is exactly $h$ sufficiently near $R$.
Here is the radial interpolation used in the construction.  Choose
$0<a<b<R$ within the interval where $h'>0$, and put
$k=h(a)/(2R)>0$.  For $a\le r<R$ we have
$kr<h(a)\le h(r)$.  Choose a smooth nondecreasing cutoff $\chi$ which
is zero for $r\le a$ and one for $r\ge b$, and set
\[
 f(r)=(1-\chi(r))kr+\chi(r)h(r)\quad(r>a),
 \qquad f(r)=kr\quad(r\le a).
\]
The cutoff makes the two definitions smooth at $a$.  On the transition
annulus,
\[
 f'(r)=(1-\chi(r))k+\chi(r)h'(r)+\chi'(r)(h(r)-kr)>0.
\]
The same positivity holds on either remaining interval.  The linear
germ at zero, divergence at $R$, and the intermediate value theorem
give bijectivity onto $(0,\infty)$; the inverse function theorem gives
a smooth inverse.

The radial diffeomorphism $x\mapsto f(\|x\|)x/\|x\|$ of
$B(0,R)$ onto $\mathbb R^3$ is linear near zero.
Compose it with $K^{-1}$ and the exterior ball chart.  After choosing
a common small $\epsilon$, the formula is
$f(\rho(s))q=(1/s)q$ on every angular direction.
\end{proof}

\begin{theorem}[Binary sphere identity for the supplied gluing]
\label{thm:binary-sphere-identity}
If the data of Definition~\ref{def:sphere-binary-data} are given and
$A,B$ are diffeomorphic to $S^3$, then $C$ is diffeomorphic to $S^3$.
\lean{PoincareMT.SmoothConnectedSumData.nonempty_diffeomorph_threeSphere}
\source{MT2007}
\dcref{Corollary 15.4(2), pp. 358--359}
\uses{def:sphere-binary-data,lem:reciprocal-sphere-end}
\end{theorem}

\begin{proof}
Choose reciprocal end coordinates $E_A,E_B$.  Apply the radial extension
of the sphere isotopy to $\varphi^{-1}$, making it equal to
$rq\mapsto r\varphi^{-1}(q)$ for $r\ge2$, and compose it with $E_B$.
The resulting Euclidean coordinates $u_A$ on $V_A$ and $u_B$ on $V_B$
satisfy, on a common collar $|s|<w$ with $0<w\le1/2$,
\[
 u_A(c(q,s))=-s^{-1}q\quad(s<0),\qquad
 u_B(c(q,s))=s^{-1}q\quad(s>0).
\]
The bound $w\le1/2$ ensures that the positive-end radius is at least two,
where the angular correction is exactly $\varphi^{-1}$.

Define the smooth radial compression and inversion
\[
 H(z)=\frac{2z}{\sqrt{1+\|z\|^2}+1},\qquad
 J(z)=\frac{4z}{\|z\|^2}\quad(z\ne0).
\]
The first is a diffeomorphism $\mathbb R^3\to B(0,2)$,
with inverse $H^{-1}(y)=4y/(4-\|y\|^2)$.
The second is a smooth involution away from zero.  Let
$o_A=u_A^{-1}(0)$ and $o_B=u_B^{-1}(0)$.  On $C\setminus\{o_B\}$
set
\[
 \Psi(x)=
 \begin{cases}
 H(u_A(x)),&x\in V_A,\\
 J(H(u_B(x))),&x\in V_B\setminus\{o_B\},\\
 2q,&x=c(q,0).
 \end{cases}
\]
These three regions map respectively to radii less than, greater than,
and equal to two.  Their inverse formulas therefore give a bijection
onto $\mathbb R^3$.

Smoothness across the central sphere is a calculation, not a consequence
of piecewise continuity.  On the whole smaller collar the three
expressions equal
\[
       \Psi(c(q,s))=a(s)q,\qquad
       a(s)=2\bigl(\sqrt{1+s^2}+s\bigr).
\]
Indeed $(\sqrt{1+s^2}+s)(\sqrt{1+s^2}-s)=1$.
The radius is positive and strictly increasing, and its inverse is
$s=r/4-1/r$.  Thus both $\Psi$ and its inverse are smooth at radius two.
They are already smooth on the cap regions by their explicit formulas.

Interchange the two caps and replace $s$ by $-s$ to obtain a second
Euclidean chart on $C\setminus\{o_A\}$.  Since $a(s)a(-s)=4$,
the transition on the overlap is $J$.  The same-basis stereographic
charts of the standard unit $S^3$ have precisely this transition
(the factor four reflects the chosen stereographic normalization).
Compose each chart on $C$ with the inverse standard chart.  The
transition identity makes the maps agree on the overlap; their inverses
agree as well.  They glue to the required global diffeomorphism.
\end{proof}

\section{Clopen components and finite induction}
\label{sec:sphere-union-induction}

\begin{definition}[Finite sphere union]
\label{def:sphere-union}
A manifold $X$ is a finite sphere union if it has a smooth disjoint-union
presentation by finitely many manifolds diffeomorphic to $S^3$.
The empty family is permitted and presents only the empty manifold.
\lean{PoincareMT.M74.SphereUnion}
\source{MT2007}
\dcref{Componentwise implementation of Corollary 15.4(2), pp. 358--359}
\end{definition}

\begin{lemma}[Preservation under one operation]
\label{lem:sphere-union-step}
A binary connected-sum operation whose source is a finite sphere union
has a target which is a finite sphere union.
\lean{PoincareMT.M74.SphereUnion.of_connectedSumStep}
\source{MT2007}
\dcref{Corollary 15.4(2), pp. 358--359}
\uses{def:sphere-union,thm:binary-sphere-identity}
\end{lemma}

\begin{proof}
Let the old source $X$ be partitioned into sphere regions $D_i$, and
also into the two clopen side regions representing $A$ and $B$ for
the operation.  Each radius-two ball chart is connected and nonempty.
A connected subset of a finite clopen partition lies in a unique
region: its intersection with one nonempty part is clopen in that
subset and hence equals it.  Thus the two ball charts select indices
$i_0,i_1$.  Each $D_i$ is itself connected; if it meets a side region,
it lies wholly in that side.  Therefore $i_0\ne i_1$, since the two
side regions are disjoint.  The selected regions contain the removed
closed balls as well as their full radius-two charts.

Pull $D_{i_0}$ back to $A$ to obtain a clopen set $U$, and pull
$D_{i_1}$ back to $B$ to obtain $V$.  In the target $Y$ form
\[
 T=a(U\setminus D_{b_A})\ \cup\
       b(V\setminus D_{b_B})\ \cup\ c(S^2\times\{0\}).
\]
This set is open: away from the central sphere it is an image of an
open punctured region, and at that sphere the full collar lies in $T$
because both ball charts lie in $U,V$.  Its complement is
\[
       a(A\setminus U)\ \cup\ b(B\setminus V),
\]
which is open since these clopen sets avoid the removed balls.
Thus $T$ is clopen.  Restrict the two ball charts and side
identifications to the selected original sphere pieces.  The
unchanged collar and the unchanged gluing map give binary
connected-sum data with target $T$.  The binary sphere identity
identifies $T$ with $S^3$.

For every $i\ne i_0,i_1$, the connected region $D_i$ lies in exactly
one side and avoids that side's ball chart.  Its image under $a$ or
$b$ is a clopen target region and is smoothly identical to the whole
old sphere piece.  To verify closedness, put the full opposite side,
the complement of $D_i$ in its own side, and the central collar
together; this is the open complementary selected region by the
previous argument.  Hence no closure point of an untouched piece
is introduced at the central sphere.

The new regions are indexed by
\[
       \{*\}\ \sqcup\ \{i:i\ne i_0,\ i\ne i_1\},
\]
where $*$ represents $T$.  Regions from different old indices are
disjoint by the old partition; images from opposite sides are
disjoint by the connected-sum data; the central sphere meets
neither side image.  These facts also give disjointness from $T$.
Every point of $Y$ lies in a side image or the central sphere.
In a side image its old index is either the selected one or an
untouched one, proving that the new regions cover $Y$.
Enumerating this finite index set gives a smooth disjoint-union
presentation by spheres.
\end{proof}

\begin{theorem}[Finite connected-sum reduction]
\label{thm:sphere-finite-reduction}
Let $(P_j)_{j<r}$ be a finite family of smooth three-manifolds, each
equipped with a diffeomorphism to $S^3$.  If a finite smooth
connected-sum assembly of this family has nonempty connected target
$X$, then $X$ is smoothly diffeomorphic to $S^3$.
\lean{PoincareMT.m74ConnectedSumReduction}
\lean{PoincareMT.M74.SphereUnion.of_reflTransGen}
\lean{PoincareMT.M74.SphereUnion.nonempty_diffeomorph_threeSphere}
\source{MT2007}
\dcref{Corollary 15.4(2), pp. 358--359}
\uses{def:sphere-union,lem:sphere-union-step}
\end{theorem}

\begin{proof}
The initial disjoint union is a finite sphere union.  Apply the
preservation lemma once for every operation in the finite sequence;
for a sequence of length zero retain the initial presentation.
This proves that $X$ is a finite sphere union.
Choose a point of $X$ and a region containing it.  That region is
nonempty and clopen in the connected $X$, so it equals $X$.
Its smooth region identification and its sphere map give the result.
The zero-factor case cannot produce a nonempty connected target by
this argument: the final region selection would have no index.
\end{proof}

\section{The smooth endpoint}
\label{sec:m74-m75-endpoint}

\begin{corollary}[Reduction of the same initial slice]
\label{cor:same-slice-sphere}
An exact finite reconstruction for a simply connected original
manifold $M$ yields a diffeomorphism $d_0:M_0\to S^3$.
\lean{PoincareMT.m74Reduction_from_M72_M73}
\source{MT2007}
\dcref{Corollary 15.4(2), pp. 358--359}
\uses{thm:exact-sphere-factors,thm:sphere-finite-reduction}
\end{corollary}

\begin{proof}
Recognize every $P_j$ by Theorem~\ref{thm:exact-sphere-factors}.
Use those maps with the supplied assembly into $M_0$, preserving its
factor indices and its target.  The finite reduction theorem gives
$d_0$ because this target is nonempty and connected.
\end{proof}

\begin{theorem}[Smooth endpoint from the constructed flow]
\label{thm:sphere-smooth-endpoint}
\uses{prop:rf-normalize,prop:foundations-projective-exclusion,prop:v4-global-extinction}
Suppose that for every compact simply connected smooth three-manifold
$M$ the preceding chapters supply a normalized metric, one controlled
surgery flow with initial identification $e_0:M\to M_0$, the local
surgery topology of that flow, and its finite extinction conclusion,
as required by the finite-history theorem.  Then every such $M$ is
smoothly diffeomorphic to $S^3$.
\lean{PoincareMT.m75EndpointInputFromExtinction}
\lean{PoincareMT.m75SmoothPoincare}
\lean{PoincareMT.m90SmoothEndpointInputs}
\source{MT2007}
\dcref{Corollary 0.2(a), p. x}
\dcref{Corollary 15.4(2), pp. 358--359}
\dcref{Theorem 15.9 and Corollary 15.10, p. 363}
\uses{cor:same-slice-sphere}
\end{theorem}

\begin{proof}
Equip $M$ with its Borel measurable structure.  Use the supplied
flow, its local topology, and its extinction conclusion to construct
the exact finite reconstruction.  The preceding corollary gives
$d_0:M_0\to S^3$.  The composite $d_0\circ e_0$ is the required
diffeomorphism.  The flow in this composite is the flow reconstructed
from extinction: neither the factor family nor the time-zero slice
is replaced.
\end{proof}

The earlier chapters' global existence and calibrated extinction
arguments construct the premise of this theorem.  The conditional
composition by itself does not establish that premise.  Likewise the
topological Poincare statement requires the later compatible-smoothing
construction and transport through its specified homeomorphism.
This chapter supplies the smooth endpoint consumed by that argument.
